\documentclass[10pt,a4paper]{amsart}
\usepackage[margin=19mm]{geometry}
\usepackage{amsmath,amssymb,mathtools}
\usepackage[scr=rsfs,cal=euler]{mathalfa}
\usepackage{xfrac}
\usepackage[unicode,hidelinks,bookmarks=false]{hyperref}
\newcommand{\ie}{i.e.\ }
\newcommand{\cf}{cf.\ }
\newcommand{\resp}{respectively\ }
\newcommand{\Subsection}{\subsection}
\providecommand{\nobreakdash}{\nobreak\hbox{-}}
\setlength{\emergencystretch}{2em}
\multlinegap0pt
\usepackage{amssymb}
\usepackage{enumerate}
\usepackage{xcolor}
\newtheorem{prop}{Proposition}
\newtheorem{lem}{Lemma}
\usepackage{cleveref}
\crefname{prop}{Proposition}{Propositions}
\crefname{lem}{Lemma}{Lemmas}
\newcommand\Ad{\operatorname{\mathsf{Ad}}}
\newcommand{\D}{\mathsf{D}}
\newcommand{\G}{\mathsf{G}}
\renewcommand{\dim}{\mathsf{dim}}
\newcommand{\disc}{\mathsf{disc}}
\renewcommand{\geq}{\geqslant}
\newcommand{\la}{\langle}
\newcommand{\lla}{\la\!\la}
\newcommand{\mg}[1]{{#1}^{\times}}
\newcommand{\ra}{\rangle}
\newcommand{\rra}{\ra\!\ra}    
\newcommand{\s}{\sigma}
\newcommand{\sq}[1]{{#1}^{\times 2}}
\newcommand{\vf}{\varphi}
\newcommand{\wt}{\widetilde}
\title{Totally decomposable algebras with involution and \texorpdfstring{$R$}{R}-triviality}
\author{M. Archita, Karim Johannes Becher}
\date{Source: arXiv:2602.10302v3 (CC BY 4.0)}
\begin{document}
\maketitle

\noindent\emph{Source and licence.} Totally decomposable algebras with involution and \texorpdfstring{$R$}{R}-triviality. Version arXiv:2602.10302v3 (CC BY 4.0), distributed by arXiv under the Creative Commons Attribution 4.0 International licence, http://creativecommons.org/licenses/by/4.0/, as recorded in the arXiv licence metadata for this exact version and shown on the abstract page https://arxiv.org/abs/2602.10302v3. Full licence notice: \texttt{licenses/arxiv-2602.10302-CC-BY-4.0.txt}.\par\smallskip

\noindent\textit{Editorial scope.} Counted unit: the article's lem ad (source0.tex lines 439--445). The article's complete original proof of it is delivered with it (source0.tex lines 446--466, one \texttt{proof} environment of 21 lines). No mathematics is written, repaired or re-derived: the statement and the proof are the article's own lines, in the article's own notation, and no retained line is substituted or re-flowed.  Retained uncounted as the source-local context the proof needs: lines 293--297; lines 322--327; lines 388--391; lines 415--419.  Nothing was omitted from any retained span, so the package carries no omission record.  No retained line contains a citation, so the wrapper carries no bibliography and no bibliography entry.  No retained line contains a figure environment, an image-inclusion command or drawing code, so no image is delivered and the package declares no asset dependency.  Editorial formatting only, in addition: an AMS \texttt{amsart} preamble that restates the article's own declarations and macros used by the retained lines, namely the environments \texttt{prop}, \texttt{lem} on separate counters, together with the standard packages those lines need.  The retained environments print in this document as Proposition 1, Lemma 1 in order of appearance, whereas the article prints the same items with the numbering of its own full document.  The complete native source of the article as distributed in the arXiv e-print for this exact version is bundled unchanged as \texttt{sources/194385/source0.tex}.
\begin{prop}\label{odddimnonzerodiv}
Let $\psi$ be an odd-dimensional quadratic form over $K$.
Then $(A,\s)$ is hyperbolic if and only if $\Ad(\psi)\otimes(A,\s)$ is hyperbolic. 
Furthermore, $$\G(\Ad(\psi)\otimes (A,\s))=\G(A,\s).$$
\end{prop}

\begin{prop}\label{simsplitreduction}
Let $Q$ be a $K$-quaternion algebra, $\nu$ its norm form and set $L=K(\nu)$.
Let $(A,\s)$ be a $K$-algebra with orthogonal involution such that $A$ is Brauer equivalent to $Q$.
Then $\s_L$ is hyperbolic if and only if $\s$ is hyperbolic.
Furthermore, $A_{L}$ is split and $\G(A,\s)=\G((A,\s)_{L})\cap \mg{K}$.
\end{prop}

\begin{prop}\label{Pfister-subform}
    Let $\vf$ and $\psi$ be Pfister forms with $\dim(\vf)>\dim(\psi)\geq 2$.
    Then $\vf_{K(\psi)}$ is hyperbolic if and only if $\psi$ is a subform of $\vf$.
\end{prop}

\begin{lem}\label{L:Pfister-fufi-G}
    Let $\alpha$ and $\rho$ be Pfister forms over $K$ such that $\rho$ is nontrivial. Let $b\in\mg{K}$. 
    Set $\psi=\alpha\otimes \lla b\rra$, $\pi=\alpha\otimes\rho$ and $\wt{\pi}=\alpha\otimes (\la b\ra\perp\rho')$.
    Then \[\G(\pi_{K(\psi)})\cap\mg{K}=\D(\pi)\cdot\D(\wt{\pi}).\]
\end{lem}

\begin{lem}\label{ad}
Let $a,b\in\mg{K}$, $Q=(a,b)_K$ and $\tau$ an orthogonal involution on $Q$ with $\disc(\tau)=a\sq{K}$.
Let $\rho$ be a Pfister form over $K$, $\psi$ an odd-dimensional quadratic form over $K$ and $$(A,\s)=\Ad(\psi\otimes\rho)\otimes (Q,\tau).$$
Let $\pi=\lla a\rra\otimes \rho$ and $\wt{\pi}=\lla a\rra\otimes (\la b\ra\perp\rho')$.
Then $$\G(A,\s)=\D(\pi)\cdot \D(\wt{\pi}).$$
    Furthermore, for any field extension $K'/K$, the form $\wt{\pi}_{K'}$ is isotropic if and only if $\s_{K'}$ is hyperbolic.
    \end{lem}

\begin{proof}
Let $\nu= \lla a,b\rra$. Set $L=K(\nu)$.
Then $Q_{L}$ and $A_{L}$ are split, and 
\[(A,\s)_{L}\simeq \Ad(\psi_L)\otimes\Ad(\rho_{L})\otimes\Ad(\lla a\rra_{L})\simeq \Ad((\psi\otimes\pi)_{L}).\]

By \Cref{simsplitreduction} and \Cref{odddimnonzerodiv}, we have
\[\G(A,\s)=\G((\psi\otimes \pi)_L)\cap \mg{K}\quad\text{ and 
}\quad\G((\psi\otimes\pi)_{L})=\G(\pi_{L}).\]
Using \Cref{L:Pfister-fufi-G}, we therefore conclude that 
\[\G(A,\s)=\G(\pi_{L})\cap\mg{K}=\D(\pi)\cdot \D(\wt{\pi}).\]

Consider now a field extension $K'/K$.
Note that 
\[\wt{\pi}_{K'(\nu)}\simeq \pi_{K'(\nu)}\qquad\text{ and }\qquad(A,\s)_{K'(\nu)}\simeq \Ad((\psi\otimes\pi)_{K'(\nu)}).\]
Assume first that $\wt{\pi}_{K'}$ is isotropic.
Then $\pi_{K'(\nu)}$ is isotropic and thus hyperbolic, because it is a Pfister form. 
It follows that $\s_{K'(\nu)}$ is hyperbolic, and hence $\s_{K'}$ is hyperbolic, by \Cref{simsplitreduction}.
Assume conversely that $\s_{K'}$ is hyperbolic.
Then $\pi_{K'(\nu)}$ is hyperbolic, and it follows by \Cref{Pfister-subform} that $\nu_{K'}$ is a subform of $\pi_{K'}$.
Since $\wt{\pi}$ is Witt equivalent to $\pi\perp-\nu$ and $\dim(\wt{\pi})=\dim(\pi)$, we conclude that $\wt{\pi}_{K'}$ is isotropic.
\end{proof}

\end{document}
